\begin{tikzpicture}[unipd/diagram,x=1cm,y=.95cm]
  % Stesse dimensioni e assi per entrambi i grafici.
  \begin{scope}[shift={(0,0)}]
    \draw[->,unipd-muted] (-1.65,0)--(1.65,0) node[right] {$x$};
    \draw[->,unipd-muted] (0,-.18)--(0,2.05) node[above] {$y$};
    \draw[very thick,unipd-primary]
      (-1.35,1.35)--(0,0)--(1.35,1.35);
    \fill[unipd-primary] (0,0) circle (1.7pt);
    \node[below left] at (0,-.12) {$0$};
    \node[align=center,font=\sffamily\footnotesize]
      at (0,-.95) {Punto angoloso\\$f(x)=|x|$};
  \end{scope}
  \begin{scope}[shift={(4.70,0)}]
    \draw[->,unipd-muted] (-1.65,0)--(1.65,0) node[right] {$x$};
    \draw[->,unipd-muted] (0,-.18)--(0,2.05) node[above] {$y$};
    \draw[very thick,unipd-primary,domain=0:1.5,samples=80]
      plot (\x,{1.3*sqrt(\x)});
    \draw[very thick,unipd-primary,domain=-1.5:0,samples=80]
      plot (\x,{1.3*sqrt(-\x)});
    \fill[unipd-primary] (0,0) circle (1.7pt);
    \node[below left] at (0,-.12) {$0$};
    \node[align=center,font=\sffamily\footnotesize]
      at (0,-.95) {Punto cuspidale\\$g(x)=\sqrt{|x|}$};
  \end{scope}
\end{tikzpicture}
